\setsubtitle{
    author : blame\\
    首殺 : 李柏霆
}
\section{F. 種馬拉車}

\begin{frame}{\btitle{題目敘述}}
    有一個長度為 $2n$ 的陣列 $A$

    上面有著 $0 \sim n-1$ 各兩個

    每個回文區間段會對答案貢獻區間最大值
\end{frame}

\begin{frame}{\btitle{前言}}
    題目叫做馬拉車，然後放在 pF

    其實根本在搞人的

    這題的定位在水題，只是問你有沒有勇氣開這題而已omob
\end{frame}

\begin{frame}{\btitle{子任務}}
    \begin{enumerate}
        \item{$1 \le n \le 100$}
        \item{$1 \le n \le 2000$}
        \item{無額外限制}
    \end{enumerate}
\end{frame}

\begin{frame}{\btitle{$1 \le n \le 100$}}
    $O(n^2)$ 枚舉每個區間

    花費 $O(n)$ 檢查是否為回文，且 $O(n)$ 計算最大值
\end{frame}

\begin{frame}{\btitle{$1 \le n \le 2000$}}
    $O(n^2)$ 枚舉每個區間

    用 rolling hash 之類的神奇演算法 $O(1)$ 檢查是否為回文

    用線段樹、sparse-table 之類的神奇算法 $O(1)$ or $O(\log n)$ 計算最大值
\end{frame}

\begin{frame}{\btitle{無額外限制}}
    官解當然不可能叫你做馬拉車

    你發現你剛剛根本沒有用到題目的性質:

    \begin{property}
        $0 \sim n-1$ 各出現兩次
    \end{property}

    而回文有性質

    \begin{property}
        對於區間(長度 $>2$) $[l, r]$ 是回文，區間 $[l+1, r-1]$ 也一定是回文
    \end{property}

    所以你可以仔細思考一下，發現這題最多也就 $O(3 n)$ 個回文而已
\end{frame}

\begin{frame}{\btitle{無額外限制}}
    作法就是由每個回文中心開始往左右擴

    假如擴增遇到左右不相同則停止擴增

    最大值也可以 $O(1)$ 維護

    所以複雜度 $O(n)$
    \pause

    假如你覺得這是 $O(n^2)$ 但寫完發現過了，那你很幸運
\end{frame}
